\import{algebra/module/kernel-image.tex}
\import{set/closure-system.tex}

\section{Spans of subsets of modules}

The span of a subset is the bounded closure hull among all submodule carriers
of a fixed ambient module.  This gives the usual least-submodule
characterization without introducing a separate finite-linear-combination
construction.

\begin{proposition}\label{submodule_carriers_form_closure_system}
    Let $R$ be a ring.
    Let $M$ be a left module over $R$.
    Let
        $C=\{B\in\pow{\carrier[M]}\mid
            \text{$B$ is a submodule of $M$ over $R$}\}$.
    Then $C$ is a closure system on $\carrier[M]$.
\end{proposition}
\begin{proof}
    Let $G=\ModuleAdditiveGroup{M}$.
    $G$ is a group by \obv{module_additive_group_is_group}.
    Show that $C\subseteq\pow{\carrier[M]}$.
    \begin{subproof}
        Show that for every $B\in C$ we have
            $B\in\pow{\carrier[M]}$.
        \begin{subproof}
            Follows by \auto.
        \end{subproof}
        Follows by \obv{subseteq}.
    \end{subproof}
    Show that for every $D\subseteq C$ we have
        $\boundedIntersection{\carrier[M]}{D}\in C$.
    \begin{subproof}
        Fix $D\subseteq C$.
        Let $P=\boundedIntersection{\carrier[M]}{D}$.
        $P\subseteq\carrier[M]$ by
            \auto{bounded_intersection_subseteq}.
        Hence $P\subseteq\carrier[G]$ by
            \auto{module_additive_group_fields}.
        Show that $\neutral[G]\in P$.
        \begin{subproof}
            $\neutral[G]\in\carrier[G]$ by
                \obv{group_is_monoid,monoid_is_unital_magma,
                    unitalmagma_type}.
            Hence $\neutral[G]\in\carrier[M]$ by
                \auto{module_additive_group_fields}.
            Show that for every $B\in D$ we have $\neutral[G]\in B$.
            \begin{subproof}
                Fix $B\in D$.
                Hence $B\in C$ by \obv{subseteq}.
                Hence $B$ is a submodule of $M$ over $R$ by
                    \auto.
                Follows by
                    \auto{submodule_zero,module_additive_group_fields}.
            \end{subproof}
            Follows by \obv{bounded_intersection_iff}.
        \end{subproof}
        Show that for all $x,y\in P$ we have $\mul[G](x,y)\in P$.
        \begin{subproof}
            Fix $x,y\in P$.
            Hence $x,y\in\carrier[G]$ by \obv{subseteq}.
            $\mul[G](x,y)\in\carrier[G]$ by
                \obv{group_is_monoid,monoid_is_magma,magma_welldef}.
            Hence $\mul[G](x,y)\in\carrier[M]$ by
                \auto{module_additive_group_fields}.
            Show that for every $B\in D$ we have $\mul[G](x,y)\in B$.
            \begin{subproof}
                Fix $B\in D$.
                Hence $B\in C$ by \obv{subseteq}.
                Hence $B$ is a submodule of $M$ over $R$ by
                    \auto.
                $x,y\in B$ by \auto{bounded_intersection_iff}.
                Follows by
                    \auto{submodule_addition,
                        module_additive_group_fields}.
            \end{subproof}
            Follows by \obv{bounded_intersection_iff}.
        \end{subproof}
        Show that for every $x\in P$ we have $\groupInv[G](x)\in P$.
        \begin{subproof}
            Fix $x\in P$.
            Hence $x\in\carrier[G]$ by \obv{subseteq}.
            $\groupInv[G](x)\in\carrier[G]$ by
                \obv{group_inverse_type}.
            Hence $\groupInv[G](x)\in\carrier[M]$ by
                \auto{module_additive_group_fields}.
            Show that for every $B\in D$ we have $\groupInv[G](x)\in B$.
            \begin{subproof}
                Fix $B\in D$.
                Hence $B\in C$ by \obv{subseteq}.
                Hence $B$ is a submodule of $M$ over $R$ by
                    \auto.
                $x\in B$ by \auto{bounded_intersection_iff}.
                Follows by
                    \auto{submodule_negation,
                        module_additive_group_fields}.
            \end{subproof}
            Follows by \obv{bounded_intersection_iff}.
        \end{subproof}
        Hence $\subgroup{P}{G}$ by \obv{subgroup}.
        Show that for all $a,x$ such that $a\in\carrier[R]$ and $x\in P$
            we have $\scalarMul[M](a,x)\in P$.
        \begin{subproof}
            Fix $a,x$.
            Assume $a\in\carrier[R]$ and $x\in P$.
            Hence $x\in\carrier[M]$ by \obv{subseteq}.
            $\scalarMul[M](a,x)\in\carrier[M]$ by
                \obv{module_scalar_type}.
            Show that for every $B\in D$ we have
                $\scalarMul[M](a,x)\in B$.
            \begin{subproof}
                Fix $B\in D$.
                Hence $B\in C$ by \obv{subseteq}.
                Hence $B$ is a submodule of $M$ over $R$ by
                    \auto.
                $x\in B$ by \auto{bounded_intersection_iff}.
                Follows by \obv{submodule_scalar}.
            \end{subproof}
            Follows by \obv{bounded_intersection_iff}.
        \end{subproof}
        Hence $P$ is a submodule of $M$ over $R$ by
            \obv{submodule}.
        $P\in\pow{\carrier[M]}$ by \obv{powerset_intro}.
        Hence $P\in C$ by \auto.
    \end{subproof}
    Follows by \obv{closure_system}.
\end{proof}

\begin{definition}\label{module_span}
    Let $R$ be a ring.
    Let $M$ be a left module over $R$.
    $\moduleSpan{R}{M}{A}=
        \closureHull{\carrier[M]}
            {\{B\in\pow{\carrier[M]}\mid
                \text{$B$ is a submodule of $M$ over $R$}\}}
            {A}$.
\end{definition}

\begin{proposition}\label{module_span_subseteq_carrier}
    Let $R$ be a ring.
    Let $M$ be a left module over $R$.
    Let $A$ be a subset of $\carrier[M]$.
    Then $\moduleSpan{R}{M}{A}\subseteq\carrier[M]$.
\end{proposition}
\begin{proof}
    Follows by \auto{module_span,closure_hull_subseteq}.
\end{proof}

\begin{proposition}\label{module_span_contains_generators}
    Let $R$ be a ring.
    Let $M$ be a left module over $R$.
    Let $A$ be a subset of $\carrier[M]$.
    Then $A\subseteq\moduleSpan{R}{M}{A}$.
\end{proposition}
\begin{proof}
    Let
        $C=\{B\in\pow{\carrier[M]}\mid
            \text{$B$ is a submodule of $M$ over $R$}\}$.
    $C$ is a closure system on $\carrier[M]$ by
        \obv{submodule_carriers_form_closure_system}.
    $A\subseteq\closureHull{\carrier[M]}{C}{A}$ by
        \obv{closure_hull_extensive}.
    Follows by \obv{module_span}.
\end{proof}

\begin{theorem}\label{module_span_is_submodule}
    Let $R$ be a ring.
    Let $M$ be a left module over $R$.
    Let $A$ be a subset of $\carrier[M]$.
    Then $\moduleSpan{R}{M}{A}$ is a submodule of $M$ over $R$.
\end{theorem}
\begin{proof}
    Let
        $C=\{B\in\pow{\carrier[M]}\mid
            \text{$B$ is a submodule of $M$ over $R$}\}$.
    $C$ is a closure system on $\carrier[M]$ by
        \obv{submodule_carriers_form_closure_system}.
    $\closureHull{\carrier[M]}{C}{A}\in C$ by
        \obv{closure_hull_closed}.
    Follows by \auto{module_span}.
\end{proof}

\begin{theorem}\label{module_span_least}
    Let $R$ be a ring.
    Let $M$ be a left module over $R$.
    Let $A$ be a subset of $\carrier[M]$.
    Suppose $B$ is a submodule of $M$ over $R$ and $A\subseteq B$.
    Then $\moduleSpan{R}{M}{A}\subseteq B$.
\end{theorem}
\begin{proof}
    Let
        $C=\{K\in\pow{\carrier[M]}\mid
            \text{$K$ is a submodule of $M$ over $R$}\}$.
    $C$ is a closure system on $\carrier[M]$ by
        \obv{submodule_carriers_form_closure_system}.
    $B\subseteq\carrier[M]$ by \obv{submodule_subseteq}.
    Hence $B\in\pow{\carrier[M]}$ by \obv{powerset_intro}.
    Hence $B\in C$.
    $\closureHull{\carrier[M]}{C}{A}\subseteq B$ by
        \obv{closure_hull_least}.
    Follows by \obv{module_span}.
\end{proof}

\begin{theorem}\label{module_span_universal_property}
    Let $R$ be a ring.
    Let $M$ be a left module over $R$.
    Let $A$ be a subset of $\carrier[M]$.
    Suppose $B$ is a submodule of $M$ over $R$.
    Then $\moduleSpan{R}{M}{A}\subseteq B$ iff $A\subseteq B$.
\end{theorem}
\begin{proof}
    $A\subseteq\moduleSpan{R}{M}{A}$ by
        \obv{module_span_contains_generators}.
    Follows by \auto{module_span_least,subseteq_transitive}.
\end{proof}

\begin{proposition}\label{module_span_monotone}
    Let $R$ be a ring.
    Let $M$ be a left module over $R$.
    Let $A,B\subseteq\carrier[M]$.
    Suppose $A\subseteq B$.
    Then $\moduleSpan{R}{M}{A}\subseteq\moduleSpan{R}{M}{B}$.
\end{proposition}
\begin{proof}
    $A\subseteq\carrier[M]$ by \obv{subseteq_transitive}.
    $\moduleSpan{R}{M}{B}$ is a submodule of $M$ over $R$ by
        \obv{module_span_is_submodule}.
    $B\subseteq\moduleSpan{R}{M}{B}$ by
        \obv{module_span_contains_generators}.
    Hence $A\subseteq\moduleSpan{R}{M}{B}$ by
        \obv{subseteq_transitive}.
    Follows by \obv{module_span_least}.
\end{proof}

\begin{corollary}\label{module_span_of_submodule}
    Let $R$ be a ring.
    Let $M$ be a left module over $R$.
    Let $A$ be a subset of $\carrier[M]$.
    Suppose $A$ is a submodule of $M$ over $R$.
    Then $\moduleSpan{R}{M}{A}=A$.
\end{corollary}
\begin{proof}
    $A\subseteq\carrier[M]$ by \obv{submodule_subseteq}.
    $A\subseteq\moduleSpan{R}{M}{A}$ by
        \obv{module_span_contains_generators}.
    $\moduleSpan{R}{M}{A}\subseteq A$ by
        \obv{module_span_least,subseteq_refl}.
    Follows by \obv{subseteq_antisymmetric}.
\end{proof}

\subsection{Elementary span laws}

\begin{theorem}\label{module_span_empty}
    Let $R$ be a ring.
    Let $M$ be a left module over $R$.
    Then $\moduleSpan{R}{M}{\emptyset}
        =\trivialSubgroupCarrier{\ModuleAdditiveGroup{M}}$.
\end{theorem}
\begin{proof}
    Let $S=\moduleSpan{R}{M}{\emptyset}$ and
        $Z=\trivialSubgroupCarrier{\ModuleAdditiveGroup{M}}$.
    $Z$ is a submodule of $M$ over $R$ by \auto{zero_submodule}.
    $\emptyset\subseteq\carrier[M]$ and $\emptyset\subseteq Z$ by
        \obv{emptyset_subseteq}.
    Hence $S\subseteq Z$ by \auto{module_span_least}.
    $S$ is a submodule of $M$ over $R$ by
        \auto{module_span_is_submodule}.
    Hence $\neutral[M]\in S$ by \obv{submodule_zero}.
    $Z=\{\neutral[\ModuleAdditiveGroup{M}]\}$ by
        \obv{trivial_subgroup_carrier}.
    Hence $Z=\{\neutral[M]\}$ by
        \auto{module_additive_group_fields}.
    Hence $Z\subseteq S$ by \auto{singleton_subset_intro}.
    Follows by \obv{subseteq_antisymmetric}.
\end{proof}

\begin{proposition}\label{module_span_singleton_contains}
    Let $R$ be a ring.
    Let $M$ be a left module over $R$.
    Let $x\in\carrier[M]$.
    Then $x\in\moduleSpan{R}{M}{\{x\}}$.
\end{proposition}
\begin{proof}
    $\{x\}\subseteq\carrier[M]$ by
        \obv{singleton_subset_intro}.
    $\{x\}\subseteq\moduleSpan{R}{M}{\{x\}}$ by
        \obv{module_span_contains_generators}.
    Follows by \auto{subseteq,singleton_iff}.
\end{proof}

\begin{theorem}\label{module_span_singleton_least_iff}
    Let $R$ be a ring.
    Let $M$ be a left module over $R$.
    Let $x\in\carrier[M]$.
    Suppose $B$ is a submodule of $M$ over $R$.
    Then $\moduleSpan{R}{M}{\{x\}}\subseteq B$ iff $x\in B$.
\end{theorem}
\begin{proof}
    $\{x\}\subseteq\carrier[M]$ by
        \obv{singleton_subset_intro}.
    $\{x\}\subseteq B$ iff $x\in B$ by
        \auto{singleton_subset_intro,singleton_subset_elim}.
    Follows by \obv{module_span_universal_property}.
\end{proof}

\begin{theorem}\label{module_span_carrier}
    Let $R$ be a ring.
    Let $M$ be a left module over $R$.
    Then $\moduleSpan{R}{M}{\carrier[M]}=\carrier[M]$.
\end{theorem}
\begin{proof}
    $\carrier[M]$ is a submodule of $M$ over $R$ by
        \obv{module_carrier_submodule}.
    Follows by
        \obv{module_span_of_submodule,subseteq_refl}.
\end{proof}

\begin{theorem}\label{module_span_union}
    Let $R$ be a ring.
    Let $M$ be a left module over $R$.
    Let $A$ be a subset of $\carrier[M]$.
    Let $B$ be a subset of $\carrier[M]$.
    Then
        $\moduleSpan{R}{M}{A\union B}
            =\moduleSpan{R}{M}{
                \moduleSpan{R}{M}{A}\union
                \moduleSpan{R}{M}{B}}$.
\end{theorem}
\begin{proof}
    Let $P=\moduleSpan{R}{M}{A}$ and
        $Q=\moduleSpan{R}{M}{B}$.
    Let $S=\moduleSpan{R}{M}{A\union B}$ and
        $T=\moduleSpan{R}{M}{P\union Q}$.
    $P\subseteq\carrier[M]$ and $Q\subseteq\carrier[M]$ by
        \auto{module_span_subseteq_carrier}.
    $A\union B\subseteq\carrier[M]$ and
        $P\union Q\subseteq\carrier[M]$ by
        \obv{union_subsets_is_subset}.
    $A\subseteq P$ and $B\subseteq Q$ by
        \obv{module_span_contains_generators}.
    Hence $A\union B\subseteq P\union Q$ by
        \obv{union_subseteq_union}.
    Hence $S\subseteq T$ by \auto{module_span_monotone}.
    $A\subseteq A\union B$ and $B\subseteq A\union B$ by
        \obv{union_upper_left,union_upper_right}.
    Hence $P\subseteq S$ and $Q\subseteq S$ by
        \auto{module_span_monotone}.
    Hence $P\union Q\subseteq S$ by \obv{subseteq_union_iff}.
    $S$ is a submodule of $M$ over $R$ by
        \auto{module_span_is_submodule}.
    Hence $T\subseteq S$ by \auto{module_span_least}.
    Follows by \obv{subseteq_antisymmetric}.
\end{proof}

\begin{theorem}\label{module_span_union_least_iff}
    Let $R$ be a ring.
    Let $M$ be a left module over $R$.
    Let $A$ be a subset of $\carrier[M]$.
    Let $B$ be a subset of $\carrier[M]$.
    Suppose $Q$ is a submodule of $M$ over $R$.
    Then $\moduleSpan{R}{M}{A\union B}\subseteq Q$ iff
        $A\subseteq Q$ and $B\subseteq Q$.
\end{theorem}
\begin{proof}
    $A\union B\subseteq\carrier[M]$ by
        \obv{union_subsets_is_subset}.
    $A\union B\subseteq Q$ iff $A\subseteq Q$ and $B\subseteq Q$ by
        \obv{subseteq_union_iff}.
    Follows by \obv{module_span_universal_property}.
\end{proof}

\subsection{Sums of submodules}

\begin{definition}\label{submodule_sum}
    Let $R$ be a ring.
    Let $M$ be a left module over $R$.
    $\submoduleSum{R}{M}{N}{P}=\moduleSpan{R}{M}{N\union P}$.
\end{definition}

\begin{theorem}\label{submodule_sum_is_submodule}
    Let $R$ be a ring.
    Let $M$ be a left module over $R$.
    Suppose $N$ is a submodule of $M$ over $R$ and
        $P$ is a submodule of $M$ over $R$.
    Then $\submoduleSum{R}{M}{N}{P}$ is a submodule of $M$ over $R$.
\end{theorem}
\begin{proof}
    $N\subseteq\carrier[M]$ and $P\subseteq\carrier[M]$ by
        \obv{submodule_subseteq}.
    Hence $N\union P\subseteq\carrier[M]$ by
        \obv{union_subsets_is_subset}.
    Follows by \auto{submodule_sum,module_span_is_submodule}.
\end{proof}

\begin{proposition}\label{submodule_sum_left}
    Let $R$ be a ring.
    Let $M$ be a left module over $R$.
    Suppose $N$ is a submodule of $M$ over $R$ and
        $P$ is a submodule of $M$ over $R$.
    Then $N\subseteq\submoduleSum{R}{M}{N}{P}$.
\end{proposition}
\begin{proof}
    $N\subseteq\carrier[M]$ and $P\subseteq\carrier[M]$ by
        \obv{submodule_subseteq}.
    $N\union P\subseteq\carrier[M]$ by
        \obv{union_subsets_is_subset}.
    $N\subseteq N\union P$ by \obv{union_upper_left}.
    $N\union P\subseteq\moduleSpan{R}{M}{N\union P}$ by
        \obv{module_span_contains_generators}.
    Follows by \obv{submodule_sum,subseteq_transitive}.
\end{proof}

\begin{proposition}\label{submodule_sum_right}
    Let $R$ be a ring.
    Let $M$ be a left module over $R$.
    Suppose $N$ is a submodule of $M$ over $R$ and
        $P$ is a submodule of $M$ over $R$.
    Then $P\subseteq\submoduleSum{R}{M}{N}{P}$.
\end{proposition}
\begin{proof}
    $N\subseteq\carrier[M]$ and $P\subseteq\carrier[M]$ by
        \obv{submodule_subseteq}.
    $N\union P\subseteq\carrier[M]$ by
        \obv{union_subsets_is_subset}.
    $P\subseteq N\union P$ by \obv{union_upper_right}.
    $N\union P\subseteq\moduleSpan{R}{M}{N\union P}$ by
        \obv{module_span_contains_generators}.
    Follows by \obv{submodule_sum,subseteq_transitive}.
\end{proof}

\begin{theorem}\label{submodule_sum_iff}
    Let $R$ be a ring.
    Let $M$ be a left module over $R$.
    Suppose $N$ is a submodule of $M$ over $R$ and
        $P$ is a submodule of $M$ over $R$.
    Then $x\in\submoduleSum{R}{M}{N}{P}$ iff
        there exists $n\in N$ such that there exists $p\in P$ such that
            $x=\mul[M](n,p)$.
\end{theorem}
\begin{proof}
    Let $T=\{\mul[M](n,p)\mid n\in N,p\in P\}$.
    $M$ is a commutative group by
        \obv{module_is_commutative_group}.
    Hence $M$ is a monoid by \obv{group_is_monoid}.
    $N,P\subseteq\carrier[M]$ by \obv{submodule_subseteq}.
    Show that $T$ is a submodule of $M$ over $R$.
    \begin{subproof}
        Show that $T\subseteq\carrier[M]$.
        \begin{subproof}
            Show that every element of $T$ is an element of $\carrier[M]$.
            \begin{subproof}
                Fix $z\in T$.
                Take $n\in N$ such that there exists $p\in P$ such that
                    $z=\mul[M](n,p)$.
                Take $p\in P$ such that $z=\mul[M](n,p)$.
                $n,p\in\carrier[M]$ by \obv{subseteq}.
                Follows by
                    \auto{monoid_is_magma,magma_welldef}.
            \end{subproof}
            Follows by \obv{subseteq}.
        \end{subproof}
        $\neutral[M]\in N$ and $\neutral[M]\in P$ by
            \obv{submodule_zero}.
        $\neutral[M]\in\carrier[M]$ by \obv{unitalmagma_type}.
        $\mul[M](\neutral[M],\neutral[M])=\neutral[M]$ by
            \obv{unitalmagma_left}.
        Hence $\neutral[M]\in T$.
        Show that for all $u,v\in T$ we have $\mul[M](u,v)\in T$.
        \begin{subproof}
            Fix $u,v\in T$.
            Take $n\in N$ such that there exists $p\in P$ such that
                $u=\mul[M](n,p)$.
            Take $p\in P$ such that $u=\mul[M](n,p)$.
            Take $m\in N$ such that there exists $q\in P$ such that
                $v=\mul[M](m,q)$.
            Take $q\in P$ such that $v=\mul[M](m,q)$.
            $n,m,p,q\in\carrier[M]$ by \obv{subseteq}.
            $\mul[M](n,p),\mul[M](m,q),
                \mul[M](n,m),\mul[M](p,q)\in\carrier[M]$ by
                \obv{monoid_is_magma,magma_welldef}.
            $\mul[M](n,m)\in N$ and $\mul[M](p,q)\in P$ by
                \obv{submodule_addition}.
            $\mul[M](\mul[M](n,p),m)=
                \mul[M](\mul[M](n,m),p)$ by
                \obv{monoid_commutative_monoid_exchange,
                    commutative_group_commutative}.
            $\mul[M](\mul[M](n,p),\mul[M](m,q))=
                \mul[M](\mul[M](\mul[M](n,p),m),q)$ by
                \obv{monoid_assoc}.
            Hence $\mul[M](\mul[M](n,p),\mul[M](m,q))=
                \mul[M](\mul[M](\mul[M](n,m),p),q)$.
            $\mul[M](\mul[M](\mul[M](n,m),p),q)=
                \mul[M](\mul[M](n,m),\mul[M](p,q))$ by
                \obv{monoid_assoc}.
            Hence $\mul[M](\mul[M](n,p),\mul[M](m,q))=
                \mul[M](\mul[M](n,m),\mul[M](p,q))$.
            Follows by \auto.
        \end{subproof}
        Show that for every $u\in T$ we have $\groupInv[M](u)\in T$.
        \begin{subproof}
            Fix $u\in T$.
            Take $n\in N$ such that there exists $p\in P$ such that
                $u=\mul[M](n,p)$.
            Take $p\in P$ such that $u=\mul[M](n,p)$.
            $n,p\in\carrier[M]$ by \obv{subseteq}.
            $\groupInv[M](n)\in N$ and $\groupInv[M](p)\in P$ by
                \obv{submodule_negation}.
            $\groupInv[M](u)=
                \mul[M](\groupInv[M](n),\groupInv[M](p))$ by
                \auto{commutative_group_inverse_product}.
            Follows by \auto.
        \end{subproof}
        Show that for all $a,u$ such that $a\in\carrier[R]$ and $u\in T$
            we have $\scalarMul[M](a,u)\in T$.
        \begin{subproof}
            Fix $a,u$.
            Assume $a\in\carrier[R]$ and $u\in T$.
            Take $n\in N$ such that there exists $p\in P$ such that
                $u=\mul[M](n,p)$.
            Take $p\in P$ such that $u=\mul[M](n,p)$.
            $n,p\in\carrier[M]$ by \obv{subseteq}.
            $\scalarMul[M](a,n)\in N$ and
                $\scalarMul[M](a,p)\in P$ by
                \obv{submodule_scalar}.
            $\scalarMul[M](a,u)=
                \mul[M](\scalarMul[M](a,n),
                    \scalarMul[M](a,p))$ by
                \auto{module_scalar_vector_distributive}.
            Follows by \auto.
        \end{subproof}
        Follows by \auto{submodule_expanded_iff}.
    \end{subproof}
    Show that $N\subseteq T$.
    \begin{subproof}
        Show that every element of $N$ is an element of $T$.
        \begin{subproof}
            Fix $n\in N$.
            $n\in\carrier[M]$ by \obv{subseteq}.
            $\neutral[M]\in P$ by \obv{submodule_zero}.
            $\mul[M](n,\neutral[M])=n$ by
                \obv{monoid_is_unital_magma,unitalmagma_right}.
        \end{subproof}
        Follows by \obv{subseteq}.
    \end{subproof}
    Show that $P\subseteq T$.
    \begin{subproof}
        Show that every element of $P$ is an element of $T$.
        \begin{subproof}
            Fix $p\in P$.
            $p\in\carrier[M]$ by \obv{subseteq}.
            $\neutral[M]\in N$ by \obv{submodule_zero}.
            $\mul[M](\neutral[M],p)=p$ by
                \obv{monoid_is_unital_magma,unitalmagma_left}.
        \end{subproof}
        Follows by \obv{subseteq}.
    \end{subproof}
    $N\union P\subseteq T$ by
        \obv{union_subsets_is_subset}.
    $N\union P\subseteq\carrier[M]$ by
        \obv{union_subsets_is_subset}.
    Hence $\moduleSpan{R}{M}{N\union P}\subseteq T$ by
        \obv{module_span_least}.
    Hence $\submoduleSum{R}{M}{N}{P}\subseteq T$ by
        \obv{submodule_sum}.
    $N\subseteq\submoduleSum{R}{M}{N}{P}$ and
        $P\subseteq\submoduleSum{R}{M}{N}{P}$ by
        \obv{submodule_sum_left,submodule_sum_right}.
    Show that if $x\in\submoduleSum{R}{M}{N}{P}$, then
        there exists $n\in N$ such that there exists $p\in P$ such that
            $x=\mul[M](n,p)$.
    \begin{subproof}
        Assume $x\in\submoduleSum{R}{M}{N}{P}$.
        Hence $x\in T$ by \obv{subseteq}.
    \end{subproof}
    Show that if there exists $n\in N$ such that
        there exists $p\in P$ such that $x=\mul[M](n,p)$, then
        $x\in\submoduleSum{R}{M}{N}{P}$.
    \begin{subproof}
        Assume there exists $n\in N$ such that
            there exists $p\in P$ such that $x=\mul[M](n,p)$.
        Take $n\in N$ such that there exists $p\in P$ such that
            $x=\mul[M](n,p)$.
        Take $p\in P$ such that $x=\mul[M](n,p)$.
        Hence $n,p\in\submoduleSum{R}{M}{N}{P}$ by \obv{subseteq}.
        $\submoduleSum{R}{M}{N}{P}$ is a submodule of $M$ over $R$ by
            \obv{submodule_sum_is_submodule}.
        Hence $\mul[M](n,p)\in\submoduleSum{R}{M}{N}{P}$ by
            \obv{submodule_addition}.
    \end{subproof}
    Follows by \auto.
\end{proof}

\begin{theorem}\label{submodule_sum_least}
    Let $R$ be a ring.
    Let $M$ be a left module over $R$.
    Suppose $N$ is a submodule of $M$ over $R$ and
        $P$ is a submodule of $M$ over $R$ and
        $Q$ is a submodule of $M$ over $R$.
    Suppose $N\subseteq Q$ and $P\subseteq Q$.
    Then $\submoduleSum{R}{M}{N}{P}\subseteq Q$.
\end{theorem}
\begin{proof}
    $N\subseteq\carrier[M]$ and $P\subseteq\carrier[M]$ by
        \obv{submodule_subseteq}.
    $N\union P\subseteq\carrier[M]$ by
        \obv{union_subsets_is_subset}.
    $N\union P\subseteq Q$ by \auto{subseteq_union_iff}.
    Follows by \auto{submodule_sum,module_span_least}.
\end{proof}

\begin{theorem}\label{submodule_sum_universal_property}
    Let $R$ be a ring.
    Let $M$ be a left module over $R$.
    Suppose $N$ is a submodule of $M$ over $R$ and
        $P$ is a submodule of $M$ over $R$ and
        $Q$ is a submodule of $M$ over $R$.
    Then $\submoduleSum{R}{M}{N}{P}\subseteq Q$ iff
        $N\subseteq Q$ and $P\subseteq Q$.
\end{theorem}
\begin{proof}
    $N\subseteq\submoduleSum{R}{M}{N}{P}$ and
        $P\subseteq\submoduleSum{R}{M}{N}{P}$ by
        \obv{submodule_sum_left,submodule_sum_right}.
    Follows by \auto{submodule_sum_least,subseteq_transitive}.
\end{proof}

\begin{theorem}\label{submodule_sum_commutative}
    Let $R$ be a ring.
    Let $M$ be a left module over $R$.
    Suppose $N$ is a submodule of $M$ over $R$ and
        $P$ is a submodule of $M$ over $R$.
    Then $\submoduleSum{R}{M}{N}{P}=\submoduleSum{R}{M}{P}{N}$.
\end{theorem}
\begin{proof}
    $N\subseteq\carrier[M]$ and $P\subseteq\carrier[M]$ by
        \obv{submodule_subseteq}.
    Follows by \auto{submodule_sum,union_comm}.
\end{proof}

\begin{theorem}\label{submodule_sum_idempotent}
    Let $R$ be a ring.
    Let $M$ be a left module over $R$.
    Suppose $N$ is a submodule of $M$ over $R$.
    Then $\submoduleSum{R}{M}{N}{N}=N$.
\end{theorem}
\begin{proof}
    $N\subseteq\carrier[M]$ by \obv{submodule_subseteq}.
    $N\union N=N$ by \obv{union_idempotent}.
    Follows by
        \auto{submodule_sum,module_span_of_submodule}.
\end{proof}

\begin{theorem}\label{submodule_sum_zero}
    Let $R$ be a ring.
    Let $M$ be a left module over $R$.
    Let $Z=\trivialSubgroupCarrier{\ModuleAdditiveGroup{M}}$.
    Suppose $N$ is a submodule of $M$ over $R$.
    Then $\submoduleSum{R}{M}{N}{Z}=N$.
\end{theorem}
\begin{proof}
    $Z$ is a submodule of $M$ over $R$ by \auto{zero_submodule}.
    $N\subseteq\submoduleSum{R}{M}{N}{Z}$ by
        \obv{submodule_sum_left}.
    $\neutral[M]\in N$ by \obv{submodule_zero}.
    $Z=\{\neutral[\ModuleAdditiveGroup{M}]\}$ by
        \obv{trivial_subgroup_carrier}.
    Hence $Z=\{\neutral[M]\}$ by
        \auto{module_additive_group_fields}.
    Hence $Z\subseteq N$ by \auto{singleton_subset_intro}.
    $\submoduleSum{R}{M}{N}{Z}\subseteq N$ by
        \obv{submodule_sum_least,subseteq_refl}.
    Follows by \obv{subseteq_antisymmetric}.
\end{proof}

\begin{theorem}\label{submodule_sum_associative}
    Let $R$ be a ring.
    Let $M$ be a left module over $R$.
    Suppose $N$ is a submodule of $M$ over $R$ and
        $P$ is a submodule of $M$ over $R$ and
        $Q$ is a submodule of $M$ over $R$.
    Then
        $\submoduleSum{R}{M}{\submoduleSum{R}{M}{N}{P}}{Q}
            =\submoduleSum{R}{M}{N}{\submoduleSum{R}{M}{P}{Q}}$.
\end{theorem}
\begin{proof}
    Let $A=\submoduleSum{R}{M}{N}{P}$ and
        $B=\submoduleSum{R}{M}{P}{Q}$.
    Let $S=\submoduleSum{R}{M}{A}{Q}$ and
        $T=\submoduleSum{R}{M}{N}{B}$.
    $A$ and $B$ are submodules of $M$ over $R$ by
        \auto{submodule_sum_is_submodule}.
    Hence $S$ and $T$ are submodules of $M$ over $R$ by
        \auto{submodule_sum_is_submodule}.
    $N\subseteq T$ and $B\subseteq T$ by
        \obv{submodule_sum_left,submodule_sum_right}.
    $P\subseteq B$ and $Q\subseteq B$ by
        \obv{submodule_sum_left,submodule_sum_right}.
    Hence $P\subseteq T$ and $Q\subseteq T$ by
        \auto{subseteq_transitive}.
    Hence $A\subseteq T$ by \auto{submodule_sum_least}.
    Hence $S\subseteq T$ by \auto{submodule_sum_least}.
    $A\subseteq S$ and $Q\subseteq S$ by
        \obv{submodule_sum_left,submodule_sum_right}.
    $N\subseteq A$ and $P\subseteq A$ by
        \auto{submodule_sum_left,submodule_sum_right}.
    Hence $N\subseteq S$ and $P\subseteq S$ by
        \auto{subseteq_transitive}.
    Hence $B\subseteq S$ by \auto{submodule_sum_least}.
    Hence $T\subseteq S$ by \auto{submodule_sum_least}.
    Follows by \auto{subseteq_antisymmetric}.
\end{proof}

\begin{theorem}\label{submodule_sum_carrier}
    Let $R$ be a ring.
    Let $M$ be a left module over $R$.
    Suppose $N$ is a submodule of $M$ over $R$.
    Then $\submoduleSum{R}{M}{N}{\carrier[M]}=\carrier[M]$.
\end{theorem}
\begin{proof}
    $\carrier[M]$ is a submodule of $M$ over $R$ by
        \obv{module_carrier_submodule}.
    Hence $\carrier[M]\subseteq
        \submoduleSum{R}{M}{N}{\carrier[M]}$ by
        \obv{submodule_sum_right}.
    $\submoduleSum{R}{M}{N}{\carrier[M]}\subseteq\carrier[M]$ by
        \auto{submodule_sum_is_submodule,submodule_subseteq}.
    Follows by \obv{subseteq_antisymmetric}.
\end{proof}

\subsection{Linear maps and spans}

\begin{theorem}\label{linear_map_image_of_span}
    Let $R$ be a ring.
    Let $M$ be a left module over $R$.
    Let $N$ be a left module over $R$.
    Suppose $f$ is a linear map from $M$ to $N$ over $R$.
    Let $A$ be a subset of $\carrier[M]$.
    Then
        $\img{f}{\moduleSpan{R}{M}{A}}
            =\moduleSpan{R}{N}{\img{f}{A}}$.
\end{theorem}
\begin{proof}
    Let $S=\moduleSpan{R}{M}{A}$ and
        $B=\img{f}{A}$.
    Let $T=\moduleSpan{R}{N}{B}$ and $J=\img{f}{S}$.
    $M$ is a left module over $R$ and
        $N$ is a left module over $R$ by
        \obv{linear_map_source_module,linear_map_target_module}.
    $f\in\funs{\carrier[M]}{\carrier[N]}$ by
        \obv{linear_map_type}.
    Hence $B\subseteq\carrier[N]$ by
        \auto{funs_image_subseteq}.
    $S$ is a submodule of $M$ over $R$ by
        \auto{module_span_is_submodule}.
    Hence $J$ is a submodule of $N$ over $R$ by
        \auto{linear_map_submodule_image}.
    $A\subseteq S$ by \obv{module_span_contains_generators}.
    Hence $B\subseteq J$ by \auto{img_subseteq}.
    Hence $T\subseteq J$ by \auto{module_span_least}.
    $T$ is a submodule of $N$ over $R$ by
        \auto{module_span_is_submodule}.
    Hence $\preimg{f}{T}$ is a submodule of $M$ over $R$ by
        \obv{linear_map_submodule_preimage}.
    $B\subseteq T$ by \obv{module_span_contains_generators}.
    $f$ is a function and $\dom{f}=\carrier[M]$ by
        \auto{linear_map_type,funs_elim}.
    Show that $A\subseteq\preimg{f}{T}$.
    \begin{subproof}
        Show that for every $x\in A$ we have $x\in\preimg{f}{T}$.
        \begin{subproof}
            Fix $x\in A$.
            $(x,f(x))\in f$ by \auto{function_apply_elim,subseteq}.
            Hence $f(x)\in B$ by \auto{img_iff}.
            Hence $f(x)\in T$ by \obv{subseteq}.
            Follows by \auto{preimg_iff}.
        \end{subproof}
        Follows by \obv{subseteq}.
    \end{subproof}
    Hence $S\subseteq\preimg{f}{T}$ by \auto{module_span_least}.
    Show that $J\subseteq T$.
    \begin{subproof}
        Show that for every $y\in J$ we have $y\in T$.
        \begin{subproof}
            Fix $y\in J$.
            Take $x\in S$ such that $(x,y)\in f$ by \auto{img_iff}.
            Hence $x\in\preimg{f}{T}$ by \obv{subseteq}.
            Take $z\in T$ such that $(x,z)\in f$ by
                \obv{preimg_iff}.
            $f$ is right-unique by \obv{function_rightunique}.
            Hence $y=z$ by \auto{rightunique}.
            Follows by \auto.
        \end{subproof}
        Follows by \obv{subseteq}.
    \end{subproof}
    Follows by \auto{subseteq_antisymmetric}.
\end{proof}

\begin{theorem}\label{linear_maps_agree_on_span}
    Let $R$ be a ring.
    Let $M$ be a left module over $R$.
    Let $N$ be a left module over $R$.
    Suppose $f$ is a linear map from $M$ to $N$ over $R$ and
        $g$ is a linear map from $M$ to $N$ over $R$.
    Let $A$ be a subset of $\carrier[M]$.
    Suppose for every $a\in A$ we have $f(a)=g(a)$.
    Then for every $x\in\moduleSpan{R}{M}{A}$ we have $f(x)=g(x)$.
\end{theorem}
\begin{proof}
    Let $E=\{x\in\carrier[M]\mid f(x)=g(x)\}$.
    Let $G=\ModuleAdditiveGroup{M}$ and
        $H=\ModuleAdditiveGroup{N}$.
    $M$ is a left module over $R$ and
        $N$ is a left module over $R$ by
        \obv{linear_map_source_module,linear_map_target_module}.
    $\groupHom{f}{G}{H}$ and $\groupHom{g}{G}{H}$ by
        \auto{linear_map_is_additive_group_hom}.
    Hence $G$ is a group and $H$ is a group by
        \obv{group_hom_source_group,group_hom_target_group}.
    $E\subseteq\carrier[G]$ by
        \auto{subseteq,module_additive_group_fields}.
    $f(\neutral[G])=\neutral[H]$ and
        $g(\neutral[G])=\neutral[H]$ by \obv{group_hom_neutral}.
    $\neutral[G]\in\carrier[G]$ by
        \obv{group_hom_source_group,group_is_monoid,
            monoid_is_unital_magma,unitalmagma_type}.
    Hence $\neutral[G]\in\carrier[M]$ by
        \auto{module_additive_group_fields}.
    Hence $\neutral[G]\in E$.
    Show that for all $x,y\in E$ we have $\mul[G](x,y)\in E$.
    \begin{subproof}
        Fix $x,y\in E$.
        Hence $x,y\in\carrier[M]$ and
            $f(x)=g(x)$ and $f(y)=g(y)$.
        Hence $x,y\in\carrier[G]$ by \obv{subseteq}.
        $\mul[G](x,y)\in\carrier[G]$ by
            \obv{group_hom_source_group,group_is_monoid,
                monoid_is_magma,magma_welldef}.
        Hence $\mul[G](x,y)\in\carrier[M]$ by
            \auto{module_additive_group_fields}.
        $f(\mul[G](x,y))=\mul[H](f(x),f(y))$ and
            $g(\mul[G](x,y))=\mul[H](g(x),g(y))$ by
            \obv{group_hom_mul}.
        Hence $f(\mul[G](x,y))=g(\mul[G](x,y))$ by \auto.
    \end{subproof}
    Show that for every $x\in E$ we have $\groupInv[G](x)\in E$.
    \begin{subproof}
        Fix $x\in E$.
        Hence $x\in\carrier[M]$ and $f(x)=g(x)$.
        Hence $x\in\carrier[G]$ by \obv{subseteq}.
        $\groupInv[G](x)\in\carrier[G]$ by
            \obv{group_inverse_type}.
        Hence $\groupInv[G](x)\in\carrier[M]$ by
            \auto{module_additive_group_fields}.
        $f(\groupInv[G](x))=\groupInv[H](f(x))$ and
            $g(\groupInv[G](x))=\groupInv[H](g(x))$ by
            \obv{group_hom_inverse}.
        Hence $f(\groupInv[G](x))=g(\groupInv[G](x))$ by
            \auto.
    \end{subproof}
    Hence $\subgroup{E}{G}$ by \obv{subgroup}.
    Show that for all $a,x$ such that $a\in\carrier[R]$ and $x\in E$
        we have $\scalarMul[M](a,x)\in E$.
    \begin{subproof}
        Fix $a,x$.
        Assume $a\in\carrier[R]$ and $x\in E$.
        Hence $x\in\carrier[M]$ and $f(x)=g(x)$.
        $\scalarMul[M](a,x)\in\carrier[M]$ by
            \obv{linear_map_scalar_ring,module_scalar_type}.
        $f(\scalarMul[M](a,x))=\scalarMul[N](a,f(x))$ and
            $g(\scalarMul[M](a,x))=\scalarMul[N](a,g(x))$ by
            \obv{linear_map_scalar}.
        Hence $f(\scalarMul[M](a,x))=g(\scalarMul[M](a,x))$ by
            \auto.
    \end{subproof}
    Hence $E$ is a submodule of $M$ over $R$ by \auto{submodule}.
    Show that $A\subseteq E$.
    \begin{subproof}
        Show that for every $a\in A$ we have $a\in E$.
        \begin{subproof}
            Fix $a\in A$.
            Hence $a\in\carrier[M]$ by \obv{subseteq}.
            $f(a)=g(a)$.
        \end{subproof}
        Follows by \obv{subseteq}.
    \end{subproof}
    Hence $\moduleSpan{R}{M}{A}\subseteq E$ by
        \obv{module_span_least}.
    Follows by \auto{subseteq}.
\end{proof}

\begin{theorem}\label{linear_map_equality_from_spanning_set}
    Let $R$ be a ring.
    Let $M$ be a left module over $R$.
    Let $N$ be a left module over $R$.
    Suppose $f$ is a linear map from $M$ to $N$ over $R$ and
        $g$ is a linear map from $M$ to $N$ over $R$.
    Let $A$ be a subset of $\carrier[M]$.
    Suppose $\moduleSpan{R}{M}{A}=\carrier[M]$.
    Suppose for every $a\in A$ we have $f(a)=g(a)$.
    Then $f=g$.
\end{theorem}
\begin{proof}
    For every $x\in\moduleSpan{R}{M}{A}$ we have $f(x)=g(x)$ by
        \obv{linear_maps_agree_on_span}.
    Hence for every $x\in\carrier[M]$ we have $f(x)=g(x)$ by
        \auto.
    Follows by \obv{linear_map_extensionality}.
\end{proof}
